\documentclass[11pt,a4paper]{article}
\usepackage[T1]{fontenc}
\usepackage[utf8]{inputenc}
\usepackage[main=italian,provide=*]{babel}
\usepackage{amsmath,amssymb}
\usepackage{geometry}
\geometry{margin=2.3cm,top=2.5cm,bottom=2.7cm}
\usepackage{booktabs}
\usepackage[table]{xcolor}
\usepackage{tcolorbox}
\tcbuselibrary{skins,breakable}
\usepackage{titlesec}
\usepackage{enumitem}
\usepackage{seqsplit}
\usepackage{needspace}
\usepackage{fancyhdr}
\usepackage{hyperref}
\hypersetup{colorlinks=true,linkcolor=Sepia,citecolor=Sepia,urlcolor=Sepia}

\definecolor{inkblue}{HTML}{1B3A5C}
\definecolor{lightblue}{HTML}{EAF1F8}
\definecolor{accent}{HTML}{B0562A}
\definecolor{softgray}{HTML}{6B6B6B}
\definecolor{okgreen}{HTML}{2E6B3E}
\definecolor{okgreenbg}{HTML}{EAF5EC}
\definecolor{warnamber}{HTML}{9C6B12}
\definecolor{warnbg}{HTML}{FBF1E0}
\definecolor{advisorpurple}{HTML}{5B3A7A}
\definecolor{advisorbg}{HTML}{F1ECF6}
\definecolor{notebar}{HTML}{9C6B12}

\titleformat{\section}
  {\normalfont\Large\bfseries\color{inkblue}}
  {\thesection}{0.6em}{}[{\color{inkblue!40}\titlerule}]
\titlespacing{\section}{0pt}{0.8em}{0.4em}
\titleformat{\subsection}
  {\normalfont\large\bfseries\color{accent}}
  {}{0em}{}
\titlespacing{\subsection}{0pt}{0.5em}{0.2em}

\pagestyle{fancy}
\fancyhf{}
\renewcommand{\headrulewidth}{0.4pt}
\fancyhead[R]{\small\color{softgray}\thepage}
\fancyfoot[C]{\small\color{softgray}Tesi triennale, Samuele Schiavi --- Relatore: Prof.\ Alessandro Chiesa}

\newtcolorbox{keyeq}[1][]{
  colback=lightblue, colframe=inkblue!70, boxrule=0.6pt,
  arc=2pt, left=8pt, right=8pt, top=4pt, bottom=4pt, #1
}
\newtcolorbox{panelbox}[1]{
  colback=white, colframe=softgray!50, boxrule=0.5pt,
  arc=2pt, left=7pt, right=7pt, top=3pt, bottom=3pt,
  fonttitle=\bfseries\color{inkblue}, title={#1}
}
\newtcolorbox{okbox}[1]{
  colback=okgreenbg, colframe=okgreen!60, boxrule=0.6pt,
  arc=2pt, left=7pt, right=7pt, top=3pt, bottom=3pt,
  fonttitle=\bfseries\color{okgreen}, title={#1}
}
\newtcolorbox{warnbox}[1]{
  colback=warnbg, colframe=warnamber!70, boxrule=0.6pt,
  arc=2pt, left=7pt, right=7pt, top=3pt, bottom=3pt,
  fonttitle=\bfseries\color{warnamber}, title={#1}
}
\newtcolorbox{advisorbox}[1]{
  colback=advisorbg, colframe=advisorpurple!60, boxrule=0.6pt,
  arc=2pt, left=7pt, right=7pt, top=4pt, bottom=4pt,
  fonttitle=\bfseries\color{advisorpurple}, title={#1}
}
\newtcolorbox{editnote}{
  colback=white, colframe=white, boxrule=0pt,
  borderline west={2.2pt}{0pt}{notebar},
  arc=0pt, left=10pt, right=4pt, top=2pt, bottom=2pt,
  fontupper=\itshape\small\color{softgray!90!black}
}
\fancyhead[L]{\small\color{softgray}Costo in gate del circuito di Trotter}

\title{\vspace{-1.5em}\color{inkblue}Costo in gate del circuito di Trotter}
\author{Note di lavoro --- Tesi triennale, Samuele Schiavi\\ Relatore: Prof.\ Alessandro Chiesa}
\date{}

\begin{document}
\sloppy
\maketitle
\thispagestyle{fancy}
\vspace{-2em}

Analisi del numero di gate del circuito per il singolo passo di Trotter
e per $N$ passi, con la decomposizione \texttt{cx}-\texttt{rz}-\texttt{cx}
al posto del gate nativo \texttt{RZZ} (vedi la sostituzione applicata a
\texttt{\seqsplit{quantum\_simulation\_dimero\_trotter\_semplificato.ipynb}}
e \texttt{\seqsplit{quantum\_simulation\_dimero\_trotter.ipynb}}).

\section{Per un singolo passo di Trotter}

Con la decomposizione attuale (\texttt{cx}-\texttt{rz}-\texttt{cx} al
posto di \texttt{RZZ}):

\begin{center}
\renewcommand{\arraystretch}{1.25}
\begin{tabular}{@{}lc@{}}
\toprule
\rowcolor{inkblue!10}
\textbf{Tipo} & \textbf{Quantità}\\
\midrule
\texttt{rz} (campo + le 3 rotazioni sostitutive di RZZ) & 5\\
\texttt{h} (Hadamard, coniugazione del blocco DM) & 4\\
\texttt{rxx}, \texttt{ryy} & 1 + 1\\
\texttt{cx} (3 coppie CNOT--Rz--CNOT) & \textbf{6}\\
\texttt{barrier} (separatore, non un gate fisico) & 1\\
\bottomrule
\end{tabular}
\end{center}

\begin{keyeq}
\centering
\textbf{Totale: 17 gate per passo} (9 a un qubit + 8 a due qubit:
$1\,R_{XX}+1\,R_{YY}+6\,CX$), profondità 15.
\end{keyeq}

\subsection*{Scala linearmente con $N$}

$$\text{gate totali} = 17N, \qquad \text{gate a 2 qubit} = 8N, \qquad \text{profondità} = 15N$$

\begin{center}
\renewcommand{\arraystretch}{1.2}
\begin{tabular}{@{}ccccc@{}}
\toprule
\rowcolor{inkblue!10}
$N$ & gate 1-qubit & gate 2-qubit & totale & profondità\\
\midrule
1 & 9 & 8 & 17 & 15\\
2 & 18 & 16 & 34 & 30\\
5 & 45 & 40 & 85 & 75\\
\textbf{20} & \textbf{180} & \textbf{160} & \textbf{340} & \textbf{300}\\
50 & 450 & 400 & 850 & 750\\
100 & 900 & 800 & 1700 & 1500\\
\bottomrule
\end{tabular}
\end{center}

La riga $N=20$ è quella riportata nel notebook (\S 8, esecuzione con
shot a $t=30$): \textbf{160 gate a due qubit}, coerente col numero che
si vede nel notebook.

\section{Alternativa più efficiente: gate nativo \texttt{RZZ}}

La decomposizione \texttt{cx}-\texttt{rz}-\texttt{cx} non è l'unica
scelta possibile: Qiskit mette a disposizione anche il gate primitivo
\texttt{RZZGate}, richiamabile con \texttt{qc.rzz(theta, 0, 1)}, che
implementa la stessa operazione unitaria
$R_{ZZ}(\theta)=e^{-i\frac\theta2 ZZ}$ in un'unica istruzione
\textbf{astratta}, senza scomporla esplicitamente in gate a un qubit e
CNOT.

\subsection*{Da dove viene: la \texttt{.definition} di Qiskit}

\texttt{RZZGate} non è un gate "magico" indipendente: porta con sé una
\texttt{.definition} interna, che Qiskit costruisce automaticamente ed
è \textbf{esattamente} la stessa decomposizione
\texttt{cx}-\texttt{rz}-\texttt{cx} usata sopra:

\begin{panelbox}{Circuito (per $\theta=0.5$, da \texttt{RZZGate(0.5).definition})}
\begin{verbatim}
q_0: ----[X]----------[X]----
           |             |
q_1: ------*----[Rz(0.5)]-*---
\end{verbatim}
\end{panelbox}

Questo conferma che le due versioni del circuito sono \textbf{identiche
a livello di operazione fisica}: cambia solo il livello di astrazione a
cui viene scritto e contato il circuito, non il calcolo che
rappresentano.

\begin{itemize}[leftmargin=1.3em,itemsep=0.2em]
\item Scrivendo \texttt{qc.rzz(theta, 0, 1)}, Qiskit tiene l'istruzione
\textbf{come blocco unico} finché non la si scompone esplicitamente
(\texttt{qc.decompose()}) o si passa per \texttt{transpile()}.
\item Solo al \texttt{transpile()} verso un \emph{basis gate set} che
non include \texttt{rzz} nativamente (es.\ $\{$cx, rz, sx, x$\}$,
tipico di molti backend a superconduttori IBM), il transpiler la
espande automaticamente nella stessa sequenza \texttt{cx}-\texttt{rz}-\texttt{cx}
--- cioè fa lui, in automatico, il lavoro che nella versione
"esplicita" del notebook viene fatto a mano nel codice Python.
\end{itemize}

\subsection*{Perché è più efficiente (nel conteggio astratto)}

Contando le operazioni \textbf{prima} di qualunque transpile, un gate
\texttt{rzz} conta come \textbf{1} operazione, non 3
(\texttt{cx}+\texttt{rz}+\texttt{cx}). Il risparmio per passo è quindi
di $3-1=2$ operazioni per ciascuna delle tre \texttt{RZZ} presenti nel
passo (una in $H_1$, due in $H_2$), cioè \textbf{6 operazioni in meno
per passo}.

\subsection*{Per un singolo passo di Trotter (versione RZZ nativa)}

\begin{center}
\renewcommand{\arraystretch}{1.25}
\begin{tabular}{@{}lc@{}}
\toprule
\rowcolor{inkblue!10}
\textbf{Tipo} & \textbf{Quantità}\\
\midrule
\texttt{rz} (solo il campo, 2 rotazioni) & 2\\
\texttt{h} (Hadamard, coniugazione del blocco DM) & 4\\
\texttt{rxx}, \texttt{ryy} & 1 + 1\\
\texttt{rzz} (le tre esponenziali $ZZ$, come blocco unico) & \textbf{3}\\
\texttt{barrier} & 1\\
\bottomrule
\end{tabular}
\end{center}

\begin{keyeq}
\centering
\textbf{Totale: 11 gate per passo} (6 a un qubit + 5 a due qubit:
$1\,R_{XX}+1\,R_{YY}+3\,R_{ZZ}$), profondità 9.
\end{keyeq}

\subsection*{Scala linearmente con $N$}

$$\text{gate totali} = 11N, \qquad \text{gate a 2 qubit} = 5N, \qquad \text{profondità} = 9N$$

\begin{center}
\renewcommand{\arraystretch}{1.2}
\begin{tabular}{@{}ccccc@{}}
\toprule
\rowcolor{inkblue!10}
$N$ & gate 1-qubit & gate 2-qubit & totale & profondità\\
\midrule
1 & 6 & 5 & 11 & 9\\
2 & 12 & 10 & 22 & 18\\
5 & 30 & 25 & 55 & 45\\
\textbf{20} & \textbf{120} & \textbf{100} & \textbf{220} & \textbf{180}\\
50 & 300 & 250 & 550 & 450\\
100 & 600 & 500 & 1100 & 900\\
\bottomrule
\end{tabular}
\end{center}

\needspace{10\baselineskip}
\subsection*{Confronto diretto fra le due versioni}

\begin{center}
\renewcommand{\arraystretch}{1.25}
\begin{tabular}{@{}lccc@{}}
\toprule
\rowcolor{inkblue!10}
\textbf{Grandezza (per passo)} & \texttt{cx-rz-cx} esplicito & \texttt{RZZ} nativo & differenza\\
\midrule
gate a 1 qubit & 9 & 6 & $-3$\\
gate a 2 qubit & 8 & 5 & $-3$\\
totale & 17 & 11 & $-6$\\
profondità & 15 & 9 & $-6$\\
\bottomrule
\end{tabular}
\end{center}

A $N=20$ ($t=30$, come nel notebook): 340 vs 220 gate totali, 160 vs
100 gate a due qubit, profondità 300 vs 180.

\subsection*{Quando conviene l'una o l'altra}

\begin{itemize}[leftmargin=1.3em,itemsep=0.3em]
\item \textbf{\texttt{RZZ} nativo (astratto)} è la scelta più efficiente
ed è quella corretta se l'obiettivo è ragionare sul circuito a livello
di algoritmo, o se il backend di destinazione supporta \texttt{RZZ}
come gate fisico nativo (es.\ alcuni processori a ioni intrappolati o
superconduttori con accoppiamento $ZZ$ diretto): in quel caso il
conteggio astratto coincide con quello reale su hardware, e usare la
forma esplicita \texttt{cx}-\texttt{rz}-\texttt{cx} aggiungerebbe solo
gate superflui che il transpiler dovrebbe poi ri-ottimizzare.
\item \textbf{\texttt{cx}-\texttt{rz}-\texttt{cx} esplicito} è la
scelta più realistica quando il backend di destinazione ha
\textbf{solo} \texttt{CNOT} come gate nativo a due qubit (il caso più
comune sui processori a superconduttori IBM): in quel caso il
conteggio astratto di \texttt{RZZ} sottostimerebbe il costo reale,
perché il transpiler dovrà comunque espanderla in
\texttt{cx}-\texttt{rz}-\texttt{cx} più avanti --- scriverla esplicita
nel notebook rende visibile fin da subito il costo che il circuito
avrà davvero una volta compilato su quel tipo di hardware.
\end{itemize}

\begin{okbox}{In sintesi}
Le due versioni non sono in competizione come "corretta vs sbagliata",
ma rappresentano due livelli di astrazione diversi della stessa
identità matematica; quale usare dipende da cosa si vuole misurare
(costo algoritmico astratto, o costo di compilazione su un hardware
con un set di gate nativi specifico).
\end{okbox}

\end{document}
